\chapter{Вступ}\section*{Перелік умовних позначень, символів, одиниць, скорочень і термінів}\addcontentsline{toc}{section}{Перелік умовних позначень, символів, одиниць, скорочень і термінів}\begin{description}\item[QGIS] Quantum Geographic Information System.\item[OSGeo] The Open Source Geospatial Foundation.\item[ODC] Open Data Cube.\item[COAH] Copernicus Open Access Hub.\item[Sentinel-2] Набір супутників місії Sentinel-2.\item[Композит] Чотиривимірний масив із каналами Red, Green, Blue та NIR.\item[Python] Високорівнева мова програмування.\end{description}На даний час карта безхмарного покриття Землі є важливим елементом в регіональній та глобальній системах, які відстежують зміни територій, спричинені природними та техногенними факторами.

Актуальність роботи

Наприклад, оцінка шкоди, заподіяної лісовими пожежами, спричиненими вирубкою лісів, виверження вулканів, повеней або дослідження в агрокультурі, це насамперед моніторинг територій, сівозміни та інше. Для цього саме й потрібно мати безхмарний композит, щоб провести дослідження.

Також такі дані будуть корисними для багатьох картографічних веб сервісів, для яких важливо мати саме останні знімки потрібних областей.

Безхмарними знімками можуть цікавитись уряди країн, наприклад для відслідковування лісового покриву, чи детектингу інших об’єктів.

Є також інші області де потрібні такі дані, це перевезення, виробництво, лісництво, наука.

Оскільки якість знімків Землі з кожним поколінням супутників стає все більш кращою, то в майбутньому це дасть ще кращий результат розробленого алгоритму з вилучення хмар.

Мета і завдання дослідження

Головною метою роботи стала побудова моделі, за допомогою якої буде здійснюватись якісне вилучення хмарних пікселів із супутникових знімків з певної області України.

Завдання дослідження:

Ознайомитись з можливостями Open Data Cube.

Вивчити та проаналізувати вхідні дані, а саме знімки та маски хмарності.

Зробити аналіз алгоритмів машинного навчання для відновлення супутникових даних.

Побудувати модель вилучення хмар.

Прибрати візуальну різницю в кольорах між сусідніми областями супутникових знімків.

Зробити оцінку точності.

Об’єкт дослідження - супутникові дані Sentinel-2 з веб-сервісу scihub copernicus.

Предмет дослідження - морфологічний оператор розширення, регресійні методи машинного навчання.

Наукова новизна одержаних результатів

Вперше було модернізовано маску хмарності, побудована власна модель вилучення хмарних пікселів та зроблене відновлення супутникових даних.

Практичне значення одержаних результатів

Практична значимість даної роботи є досить висока, оскільки є компанії або люди, які потребують якісних безхмарних зображень територій.

Такою аудиторією, якій буде корисна така технологія може бути:

Веб-сервіси, що працюють з географічними та супутниковими даними. Наприклад сервіс моніторингу зеленого покриву.

Агро культурному сектору також потрібні дані для дослідження свого регіону.

Лісництво може використовувати дані для визначення де потрібно садити нові дерева або моніторингу області.

Державні організації, в яких є необхідність моніторингу та аналізу територій міст, наприклад урбанізація та розбудова районів, організацією якій знадобляться ці дані є підрозділ інфраструктури. Наприклад Міністерству внутрішніх справ буде корисно моніторити території для виявлення незаконного видобутку бурштину.

Міжнародні організації. Відома всім проблема скупчення великої кількості сміття в океанах або танення льодовиків. Безхмарні знімки допоможуть аналізувати ситуацію на прикладі таких проблем.

Водоканальні служби. Вони будуть спроможні аналізувати збалансоване керування водною системою певного регіону

